\chapter{Methodology}
\label{ch:methodology}

This chapter sets out the methodological spine of the thesis. It is deliberately the
longest and most detailed chapter, because the project is graded primarily on the rigour
of its evaluation rather than on orchestration sophistication. We first state the central
contribution --- the \emph{\mcp-as-time-machine} design that makes look-ahead bias
structurally impossible (\cref{sec:timemachine}) --- and the test-driven discipline that
pins it (\cref{sec:tdd}). We then treat the subtler leakage channels that a temporal
filter alone does not close: adjusted-price leakage and weight-memory contamination
(\cref{sec:leakage}). We define the market-regime labelling that segments every downstream
metric (\cref{sec:regime}), the baselines and their costing (\cref{sec:baselines}), and the
statistical treatment of uncertainty (\cref{sec:bootstrap}). The reasoning metrics
themselves are deferred to \cref{ch:evaluation}.

\section{The cardinal sin and the design response}
\label{sec:cardinal}

The dominant failure mode in LLM-driven backtesting is look-ahead bias: the agent, or the
harness around it, is exposed to information that would not have been available at the
moment of the simulated decision. KellyBench~\cite{kellybench2026} catalogues several
concrete instances --- most strikingly, label leakage introduced by fitting a calibrated
classifier on the full training set before the temporal split --- and shows that such
leakage is easy to introduce and hard to detect after the fact. The consensus across the
backtesting literature is that the only robust defence is \emph{architectural}: leakage
must be made impossible by construction, not merely discouraged by careful coding.

We take this seriously to the point of treating it as the organising principle of the
entire data layer. The design goal is a single, auditable choke point through which all
time-indexed data must pass, and a guarantee --- expressed as an executable invariant ---
that nothing dated after the simulated present can cross it.

\section{MCP-as-time-machine}
\label{sec:timemachine}

\subsection{The construction}

We introduce a simulation clock that exposes a single timestamp, \tnow. The clock is
implemented in \file{clock.py} as a small, heavily-tested module --- the one component of
the codebase the supervisor is asked to review line by line. Every \mcp\ data tool, and
every wrapper around an external data server, consults \code{get\_clock().t\_now} before
returning any datum and filters its output so that no item dated strictly after \tnow\ can
be returned. The agent calls the \emph{same} \mcp\ tools in backtest and in live operation;
only the source that advances the clock differs (an event loop in backtest, the wall clock
in live). \Cref{fig:loop} shows the resulting per-bar loop with the clock at its centre.

% ---------------------------------------------------------------------
% Figure: the perceive -> reason -> act -> observe loop with the
% simulation clock t_now at the centre, gating every tool read.
% ---------------------------------------------------------------------
\begin{figure}[t]
  \centering
  \begin{tikzpicture}[
      >=Stealth, node distance=20mm,
      stage/.style={draw, rounded corners, align=center, minimum width=26mm,
                    minimum height=11mm, fill=accent!8, draw=accent!60, thick},
      clk/.style={draw, circle, align=center, minimum size=20mm, fill=todored!8,
                  draw=todored!60, very thick},
      lbl/.style={font=\footnotesize\itshape, text=gray}
    ]
    \node[clk] (clock) {\textbf{\tnow}\\\scriptsize clock};

    \node[stage, above=18mm of clock]            (perceive) {\textbf{Perceive}\\\scriptsize read \mcp\ tools};
    \node[stage, right=24mm of clock]            (reason)   {\textbf{Reason}\\\scriptsize \CoT\ over outputs};
    \node[stage, below=18mm of clock]            (act)      {\textbf{Act}\\\scriptsize place\_order};
    \node[stage, left=24mm of clock]             (observe)  {\textbf{Observe}\\\scriptsize fills, \nav, PnL};

    % outer cycle
    \draw[->, thick] (perceive) to[bend left=18] (reason);
    \draw[->, thick] (reason)   to[bend left=18] (act);
    \draw[->, thick] (act)      to[bend left=18] (observe);
    \draw[->, thick] (observe)  to[bend left=18] (perceive);

    % every tool read consults the clock
    \draw[->, dashed, todored!70] (clock) -- (perceive) node[midway, lbl, right]{gate};
    \draw[->, dashed, todored!70] (clock) -- (act);
    \draw[->, dashed, todored!70] (clock) -- (reason);
    \draw[->, dashed, todored!70] (clock) -- (observe);

    \node[lbl, below=2mm of act, text width=58mm, align=center]
      {after the loop closes, the engine calls \code{clock.advance\_to($t{+}1$)};
       it never moves backward};
  \end{tikzpicture}
  \caption[The clock-centred agent loop]{The \mbox{perceive\,$\rightarrow$\,reason\,$\rightarrow$\,act\,$\rightarrow$\,observe}
  loop executed once per bar per ticker. The simulation clock \tnow\ sits at the centre:
  every data read (dashed arrows) is gated by it, so the agent can never perceive a
  datapoint dated after \tnow. The same loop runs in backtest and in live operation; only
  the clock's advance source differs (event loop versus wall clock). After the loop closes,
  the engine advances the clock strictly forward, which is the only mutation that exposes
  new data.}
  \label{fig:loop}
\end{figure}


This is the thesis's central methodological claim, and it is worth stating precisely.

\begin{invariant}[No look-ahead]
\label{inv:nolookahead}
For any data tool $T$ invoked while the clock reads \tnow, every item $x$ in $T$'s output
satisfies $\textsf{timestamp}(x) \le \tnow$. The bar timestamped exactly \tnow\ is
included; the first future bar is excluded. The inequality is \emph{strict} on the future
side.
\end{invariant}

The boundary semantics matter. A daily bar dated \tnow\ represents information known at the
close of that trading day, which is precisely the agent's decision point; it must be
visible. The next bar must not be. \Cref{fig:clockflow} illustrates the read path and the
strict boundary.

% ---------------------------------------------------------------------
% Figure: clock-gated data flow. raw cache -> t_now filter -> tool ->
% agent, contrasting an included past bar with an excluded future bar.
% ---------------------------------------------------------------------
\begin{figure}[t]
  \centering
  \begin{tikzpicture}[
      >=Stealth, font=\footnotesize, node distance=14mm,
      box/.style={draw, rounded corners, align=center, minimum height=10mm,
                  minimum width=24mm, thick},
      bar/.style={draw, minimum width=7mm, minimum height=6mm, font=\scriptsize},
      pass/.style={fill=accent!12, draw=accent!55},
      drop/.style={fill=todored!12, draw=todored!55},
    ]
    \node[box, fill=gray!8] (cache) {Raw parquet\\cache\\\scriptsize (unfiltered)};
    \node[box, fill=todored!8, draw=todored!60, right=18mm of cache] (filter)
      {\code{filter\_rows}\\\scriptsize keep ts $\le\tnow$};
    \node[box, fill=accent!8, right=18mm of filter] (tool) {\mcp\ tool\\\scriptsize output};
    \node[box, fill=gray!8, right=16mm of tool] (agent) {Agent};

    \draw[->, thick] (cache) -- (filter);
    \draw[->, thick] (filter) -- (tool);
    \draw[->, thick] (tool) -- (agent);

    % illustrative bars above the cache
    \node[bar, pass, above=10mm of cache.north west, anchor=west] (b1) {$t{-}2$};
    \node[bar, pass, right=1mm of b1] (b2) {$t{-}1$};
    \node[bar, pass, right=1mm of b2] (b3) {$\tnow$};
    \node[bar, drop, right=1mm of b3] (b4) {$t{+}1$};
    \node[bar, drop, right=1mm of b4] (b5) {$t{+}2$};
    \node[font=\scriptsize, above=1mm of b3] {boundary};

    \draw[->, gray] (b3.south) to[bend right=10] (cache.north);
    \draw[->, todored!70, dashed] (b4.south) to[bend left=8] (filter.north)
      node[midway, right, font=\scriptsize\itshape, text=todored]{dropped};

    \node[font=\scriptsize\itshape, text=gray, below=2mm of filter, text width=42mm,
          align=center]
      {strict inequality: the bar dated exactly \tnow\ is \emph{included}; $t{+}1$ is
       \emph{excluded}};
  \end{tikzpicture}
  \caption[Clock-gated data flow]{The clock-gated read path. Data is cached raw and never
  pre-trimmed; the \tnow\ filter is applied at read time inside the data wrapper. The
  invariant is a strict inequality --- the bar timestamped exactly \tnow\ is returned, the
  first future bar ($t{+}1$) is dropped --- so the agent sees ``the world as of the close of
  \tnow'' and nothing later. The single-item variant \code{assert\_not\_future} \emph{raises}
  instead of dropping, for cases where a future timestamp is unambiguously a bug.}
  \label{fig:clockflow}
\end{figure}


\subsection{Two enforcement modes}

The clock exposes two complementary guards, reflecting two genuinely different situations:

\begin{itemize}
  \item \code{filter\_rows(rows, date\_key)} --- the \emph{batch} guard. It returns only
  the rows with timestamp $\le\tnow$ and silently drops the rest. This is the expected,
  routine operation when slicing a cached price series to the visible window: future rows
  in the cache are not a bug, they are simply not yet visible. Dropping them is the correct
  behaviour.
  \item \code{assert\_not\_future(ts, label)} --- the \emph{single-item} guard. It
  \emph{raises} \code{FutureDataError} if \code{ts} is strictly after \tnow. This is used
  where a future timestamp is unambiguously a defect rather than an expected trim --- for
  instance, when an execution fill is stamped with a date the engine should never have
  reached.
\end{itemize}

Both guards route through one internal comparison, and both fail loudly rather than
returning a silent default. This reflects a project-wide convention: a silent fallback that
masks a leak or a missing tool call is the worst possible bug in a thesis whose entire
credibility rests on the absence of look-ahead. The codebase therefore raises on unexpected
states everywhere in the data path.

\subsection{Caching discipline that preserves the guarantee}

Performance requires caching market data to disk (parquet) so that a backtest is
re-runnable offline at zero API cost. Naively, caching and the temporal guarantee are in
tension: a cache populated by one run could expose future data to another. We resolve this
by \emph{caching raw data and applying the \tnow\ filter at read time}, never at write time.
A single parquet file is therefore valid for a run at any clock position: the same AAPL
cache serves a decision at \tnow${}={}$2023-01-03 and one at 2023-01-04, because the filter,
not the file, defines what is visible. This is shown as the substrate in
\cref{fig:architecture}.

\subsection{Timestamp normalisation}

A subtle correctness hazard is timezone handling. Comparing a timezone-aware news timestamp
against a naive \tnow\ can misplace items across the boundary by up to a day. We normalise
all timestamps to a single internal convention --- UTC-naive --- and, crucially, convert
aware datetimes to UTC \emph{before} stripping the offset. A bare offset-strip would, for
example, treat a Tokyo \texttt{+09:00} item as nine hours later than its true UTC instant,
which can flip a past item into an apparent future one and trigger a spurious leak (or, in a
filter context, an erroneous drop). The conversion is exercised in both boundary directions
by dedicated tests.

\section{Test-driven anti-leakage and mutation testing}
\label{sec:tdd}

The invariant in \cref{inv:nolookahead} is only as credible as the tests that enforce it.
Following the project's TDD mandate, the anti-look-ahead tests were written \emph{before}
the data layer they guard: \file{tests/test\_clock\_no\_lookahead.py} and
\file{tests/test\_data\_temporal\_filter.py}. They assert, among other properties, that a
tool returning a bar dated after \tnow\ causes a test failure --- which means the test
suite is red until the filter is correctly in place. These tests are treated as
non-negotiable: they are never weakened to make code pass.

Passing a hand-written test, however, only shows that the guard works on the cases the
author thought to write. To argue that the guard is \emph{load-bearing} --- that removing or
weakening it would actually be caught --- we use mutation testing as a falsification check.
Mutating the boundary comparison (for example, relaxing the strict future-side inequality
\code{>} to \code{>=}, or inverting the filter predicate) must cause at least one test to
fail. A surviving mutant would indicate a guard that is not actually exercised, which for
this particular invariant is a methodological red flag rather than a mere coverage gap. The
strict-versus-non-strict boundary is explicitly pinned this way.

\section{Leakage channels a temporal filter does not close}
\label{sec:leakage}

\Cref{inv:nolookahead} closes the \emph{when} of leakage. Two further channels concern the
\emph{what} and were the source of real, diagnosed bugs in this project; they are treated
honestly here rather than glossed over.

\subsection{Adjusted-price leakage}

Split- and dividend-adjusted price histories encode \emph{future} corporate actions into
past prices, because the adjustment factors are computed from the current state of the
world. A backtest set in 2022 that reads 2022 prices adjusted for a 2023 split is, subtly,
looking ahead. We therefore source \emph{unadjusted} prices throughout
(\code{auto\_adjust=False}) and work with raw OHLCV. The cost is that long-horizon return
comparisons are not split-corrected; this is the strictly preferable trade-off, because the
alternative is silent contamination of exactly the kind the thesis claims to prevent. The
choice is documented in the decision log.

A related, harder lesson from this project is that temporal correctness does not imply
\emph{value} correctness. An early run was corrupted not by look-ahead but by a data-handling
defect: a single-ticker download returned a column structure that, when indexed naively,
yielded a multi-element series where a scalar close price was expected, and a permissive
exception handler silently accepted malformed bars. The anti-look-ahead fortress checked
\emph{when} a bar was dated but never \emph{what} it contained. The response was a
fail-loud data-sanity gate (\file{bar\_validation.py}) that rejects non-scalar, non-finite,
non-positive, or internally inconsistent bars (e.g. high $<$ low, or a close outside a
plausible band around the open), together with a price-continuity guard that raises on
adjacent-bar jumps beyond a fixed ratio. These checks run on both the write and read paths,
so a previously cached corrupt bar is caught on first access. This episode is revisited as
part of the rigour narrative in \cref{sec:rigour}.

\subsection{Weight-memory contamination}

The most intractable channel is that the model's parameters already encode the approximate
outcome of historical events. A model trained on data through early 2025 ``knows,'' in a
diffuse statistical sense, how 2022--2024 unfolded; it can in principle infer future prices
from recalled news even when no future datum is served. KellyBench frames this directly: a
real news corpus does not eliminate the risk --- it only makes the grounding metric
meaningful, by giving cited claims something verifiable to check against.

We mitigate, and otherwise document honestly, rather than claim to solve:
\begin{enumerate}
  \item the agent is instructed to follow a rule-based process over tool outputs, reducing
  reliance on free recall;
  \item where feasible the test window is pushed toward, or past, the model's training
  cutoff, so that recall is less informative;
  \item the limitation is stated plainly as a threat to validity in \cref{sec:limitations},
  in the spirit of KellyBench's own candour.
\end{enumerate}
This honesty is itself a graded virtue: a thesis that names the channel it cannot fully
close is more credible than one that pretends the channel does not exist.

\section{Market-regime labelling}
\label{sec:regime}

Every downstream metric --- financial and reasoning alike --- is segmented by market
regime, because the thesis's substantive question is not ``can the agent trade?'' but
``\emph{when}, and \emph{why}, does it trade well or badly?''. This requires a regime
labelling that is itself causal: a label at bar $t$ must depend only on data up to $t$.

\subsection{Why rule-based, not HMM}

A hidden Markov model would give probabilistic regime assignments but requires fitting
transition parameters on historical data; done naively on the full series, that fit leaks
future information into past labels. We use instead a rule-based labelling from realised
volatility and trend, with thresholds fixed before seeing the data. It is less elegant than
an HMM but fully causal at every bar, which is the property the thesis needs; the HMM is
noted as future work.

\subsection{The detector}

The labelling (\file{mcp\_servers/analytics/regime.py}) assigns one of four labels with the
priority order \textsc{high\_vol} $>$ \textsc{bull} $>$ \textsc{bear} $>$ \textsc{range}.
The current detector (v2) improves on an initial version (v1) in two ways that are both
causality fixes:
\begin{itemize}
  \item \textbf{Reactive trend.} The bull/bear signal uses a short, 20-trading-day trend
  window ($\pm 2\%$ threshold), roughly one calendar month. The earlier 60-day window
  systematically lagged the 2022--2023 recovery, keeping a stale bearish label for weeks
  after prices had turned. The shorter window pivots about a month earlier at recovery
  inflections.
  \item \textbf{Causal volatility threshold.} The \textsc{high\_vol} label fires when
  realised volatility exceeds a trailing 252-bar percentile. The earlier version used a
  full-series percentile, which is a look-ahead: a future volatility spike would raise the
  threshold and silently re-label past bars when the series was extended. The trailing
  percentile uses only bars up to $t$. A dedicated test fails if a full-series percentile is
  reintroduced.
\end{itemize}

\subsection{Smoothing and warm-up}

Raw labels flip too often to support meaningful segmentation: a large fraction of regime
runs would otherwise last a single bar. We apply a causal minimum-hold smoothing (default
three bars, one trading week): a regime change is confirmed only after the new label has
persisted for the hold length. The smoothing is one-sided in time --- the smoothed label at
$t$ depends only on raw labels up to $t$ --- and a test verifies this anti-look-ahead
property explicitly. The raw label is retained as a diagnostic field.

The detector also imposes a hard warm-up requirement, derived rather than guessed. The
controlling constraint is the volatility-percentile window: for bar $t$ to have a
fully-populated trailing window, $t$ must be at least
$\textsf{vol\_window} + \textsf{vol\_percentile\_window} - 1 = 20 + 252 - 1 = 271$ trading
days into the data. The data window for any run therefore extends the decision window
backward by \code{REGIME\_WARMUP\_TRADING\_DAYS}${}=271$ trading days
(\code{REGIME\_WARMUP\_CALENDAR\_DAYS}${}=420$ calendar days with a safety margin). Both
backtest engines import this single constant, and an anti-regression test fails if it is set
too small or if the percentile window is changed without realigning it --- so that the
warm-up cannot silently drift out of sync with the detector. This guarantees that the first
\emph{decision} bar already has a correctly-initialised regime label, and that no decision
is excluded as warm-up.

\section{Baselines}
\label{sec:baselines}

An agent's performance is meaningless without a systematic yardstick, and the yardstick must
be at least as carefully constructed as the agent. We implement three baselines
(\file{backtest/baselines.py}): equal-weight Buy-\&-Hold, time-series momentum, and
Bollinger-band mean-reversion. Two properties make the comparison fair.

\paragraph{One-bar signal shift.} Every baseline signal is computed from the close of bar
$t$ and executed at bar $t{+}1$ (\code{exec\_signal = signal.shift(1)}). Same-bar execution
would let a baseline trade on a close it could not yet have known --- the baselines' own
look-ahead. A test fails if the shift is removed, mirroring the structural discipline applied
to the agent.

\paragraph{Symmetric transaction costs.} Costs were, in an earlier iteration, asymmetric ---
the agent paid nothing while the baselines paid a round-trip friction --- which silently
invalidated any comparison. We now charge $5$\,bps commission plus $5$\,bps slippage
($10$\,bps one-way) to \emph{both} the agent and the baselines, with Buy-\&-Hold charged an
entry and an exit cost. A test asserts cost symmetry on every run. We additionally report
turnover and a \emph{cost drag} (total commission as a fraction of initial \nav, in bps), so
that a high-turnover strategy's frictional cost is visible rather than buried in the net
return.

The baselines are computed vectorised in pandas, which is sufficient for daily data at this
scale; the event-driven engine is reserved for the agent path, where bar-by-bar drip-feed
realism matters. The decision to use pandas rather than a heavier vectorised backtester is
recorded in the decision log and turns on readability, not performance.

\section{Quantifying uncertainty: bootstrap confidence intervals}
\label{sec:bootstrap}

A single Sharpe number from a single run proves little; KellyBench shows single-seed
estimates can vary widely. We therefore report a bootstrap confidence interval around the
Sharpe ratio (\code{bootstrap\_sharpe\_ci} in \file{eval/financial.py}), resampling daily
returns to form the interval. Two honest caveats apply.

First, we use an \emph{i.i.d.}\ bootstrap rather than the stationary bootstrap of
Politis and Romano~\cite{politis1994}. Daily equity returns have weak autocorrelation, so
the i.i.d.\ resample is a reasonable approximation, but for a strategy with genuine momentum
the interval width may be mildly understated. The stationary bootstrap is the principled
upgrade and is flagged as such.

Second, the interval is only meaningful with enough \nav\ points: on a very short window the
resampled distribution degenerates and the interval collapses to a point. We treat
\minNavPoints\ daily \nav\ points as the minimum for a reportable interval, and we
\emph{never} annualise a Sharpe ratio computed on a short window --- annualising a Sharpe
estimated over, say, nineteen trading days inflates a noisy monthly figure into a
meaningless headline number. The full-period run is what supplies the thesis's confidence
interval; short windows are used only for plumbing validation and are labelled as such.

\section{Run scale tiers and reproducibility}
\label{sec:tiers}

LLM experiments cost real money, so every run declares a \emph{tier} (\file{backtest/tier.py})
that bounds its size before any model is called: \textsc{smoke} (development and CI, no cost
acknowledgement), \textsc{medium} (first real per-regime validation, up to about one month),
and \textsc{full} (explicit thesis-scale runs). All tiers enforce the regime warm-up. Every
result in this thesis is labelled with its scope: the older two-year single-agent run supports
the primary reasoning diagnosis, while the one-year common-reference run supplies the clean
cost-symmetric comparison. Development uses the cheapest capable model with cached responses
keyed by prompt hash, so debugging re-runs cost nothing and thesis-scale runs are executed only
after dry-run guardrails and data-cache preflight pass.
